\documentclass[10pt,a4paper]{amsart}
\usepackage[margin=19mm]{geometry}
\usepackage{amsmath,amssymb,mathtools}
\usepackage[scr=rsfs,cal=euler]{mathalfa}
\usepackage{xfrac}
\usepackage[unicode,hidelinks,bookmarks=false]{hyperref}
\newcommand{\ie}{i.e.\ }
\newcommand{\cf}{cf.\ }
\newcommand{\resp}{respectively\ }
\newcommand{\Subsection}{\subsection}
\providecommand{\nobreakdash}{\nobreak\hbox{-}}
\setlength{\emergencystretch}{2em}
\multlinegap0pt
\usepackage{bm}
\usepackage{bbm}
\usepackage{dsfont}
\usepackage{xspace}
\usepackage{cancel}
\usepackage{mathtools}
\usepackage{amsthm}
\makeatletter
\@ifundefined{theorem}{\newtheorem{theorem}{Theorem}}{}
\@ifundefined{lemma}{\newtheorem{lemma}[theorem]{Lemma}}{}
\@ifundefined{proposition}{\newtheorem{proposition}[theorem]{Proposition}}{}
\makeatother
\newcommand{\C}{\mathbb C}
\newcommand{\R}{\mathbb R}
\newcommand{\ip}[2]{\left\langle #1,#2\right\rangle}
\newcommand{\spn}{\operatorname{span}}
\title[A Quantum Latin Square of Order Six with Cardinality 29]{A Quantum Latin Square of Order Six with Cardinality 29}
\author{Aishwarya P. Das, Durgesh Kumar}
\date{Source: arXiv:2608.12607v1 (CC BY 4.0)}
\begin{document}
\maketitle

\noindent\emph{Source and licence.} A Quantum Latin Square of Order Six with Cardinality 29. Version arXiv:2608.12607v1 (CC BY 4.0), distributed under the Creative Commons Attribution 4.0 International licence, as recorded in the arXiv licence metadata for this exact version and shown on the abstract page https://arxiv.org/abs/2608.12607v1. Full rights notice: licenses/arxiv-2608.12607-CC-BY-4.0.txt.\par\smallskip

\noindent\textit{Editorial scope.} Counted unit: the Proposition whose source label is \texttt{prop:V29} in the article (source0.tex lines 430--432, source label \texttt{prop:V29}), delivered together with the article's own complete original proof of it (source0.tex lines 434--510, one \texttt{proof} environment of 77 lines). No mathematics is written, repaired or re-derived: statement and proof are the article's own text line for line. Required context retained uncounted: lem:missing-column (lines 183--188); eq:completion (lines 340--348); eq:v4041 (lines 367--372); eq:aabb (lines 385--390); eq:compatibility (lines 393--395); eq:v5051 (lines 417--421). Every definition, display and label of those lines is retained unchanged. Omitted from this package and not redistributed: 1--182, 189--339, 349--366, 373--384, 391--392, 396--416, 422--429, 433--433, 511--762. Nothing inside a retained excerpt is omitted or altered: every retained body is delivered exactly as the article's lines are written, so no omission receipt of any kind accompanies this package. Citations: the retained excerpts carry no citation command at all, so no bibliography entry of the article is transcribed and no reference is redistributed. Reference adaptation: none was needed and none was made; every reference inside the delivered unit resolves inside it, so no unresolved reference or citation is produced. Printed slips kept literal: no printed word, symbol, hypothesis or number of the counted statement or of its proof was altered. Layout and class: the article's own class is \texttt{article} with packages including geometry, amsmath, amssymb, amsthm, mathtools, booktabs, microtype, hyperref, lmodern; the wrapper is the lane's pinned \texttt{amsart} editorial wrapper at 10pt on A4, which restates the theorem-like environments the retained text needs and adds the article's own macro definitions that the retained text needs (\texttt{\textbackslash C}, \texttt{\textbackslash R}, \texttt{\textbackslash ip}, \texttt{\textbackslash spn}), with extra wrapper packages (bm, bbm, dsfont, xspace, cancel, mathtools). The delivered numbering is the unit's own. Assets: no retained span contains a figure environment, a graphics-inclusion command, a PDF-page inclusion, a table or a TikZ picture; no image file is copied into this package, the package declares no asset dependency and no image file is redistributed with it. Licence: version arXiv:2608.12607v1 of this article is distributed under the Creative Commons Attribution 4.0 International licence, as recorded in the arXiv licence metadata for this exact version and shown on the abstract page https://arxiv.org/abs/2608.12607v1. A full rights notice is delivered at licenses/arxiv-2608.12607-CC-BY-4.0.txt. Characters: every retained excerpt is pure ASCII, so the delivered engine, XeLaTeX with its default Latin Modern text font, reports no missing character.
\begin{lemma}[One missing column]\label{lem:missing-column}
Let $V=(v_{ij})_{i\ne j}$ be an order-$n$ array of unit vectors in
$\C^{n-1}$.  If every punctured row and all but one punctured column are
orthonormal bases, then the remaining punctured column is also an
orthonormal basis.
\end{lemma}

\begin{equation}\label{eq:completion}
\mathcal C_z(x,y)=
\left(
\frac{\ip{z}{y}x-\ip{z}{x}y}{\rho},
\frac{\ip{z}{x}x+\ip{z}{y}y}{\rho}
\right),
\qquad
\rho=\sqrt{\ip{z}{x}^2+\ip{z}{y}^2}.
\end{equation}

\begin{equation}\label{eq:v4041}
\tau=\ip{e_0}{T},\qquad
r_{40}=\sqrt{\alpha^2+\tau^2},\qquad
v_{40}=\frac{\alpha Y+\tau T}{r_{40}},\qquad
v_{41}=\frac{\tau Y-\alpha T}{r_{40}}.
\end{equation}

\begin{equation}\label{eq:aabb}
a_1=\ip{v_{20}}{v_{13}},\quad
a_2=\ip{v_{20}}{v_{43}},\quad
b_1=\ip{v_{30}}{v_{13}},\quad
b_2=\ip{v_{30}}{v_{43}}.
\end{equation}

\begin{equation}\label{eq:compatibility}
a_1b_2-a_2b_1=0.
\end{equation}

\begin{equation}\label{eq:v5051}
r_{50}=\sqrt{a_1^2+a_2^2},\qquad
v_{50}=\frac{a_2v_{13}-a_1v_{43}}{r_{50}},\qquad
v_{51}=\frac{a_1v_{13}+a_2v_{43}}{r_{50}}.
\end{equation}

\begin{proposition}\label{prop:V29}
The array $V_{29}$ is a punctured orthonormal array of order six in $\R^5$.
\end{proposition}

\begin{proof}
We verify the punctured rows first.  We then verify columns $0,2,3,4,5$ and
use Lemma~\ref{lem:missing-column} for the remaining column.

\emph{Rows.}
Row~$0$ is the standard basis.  In row~$1$, $Y$ is perpendicular
to $\spn\{W,e_4,U,e_1\}$.  This four-dimensional space is the orthogonal
direct sum of $\spn\{v_{12},v_{15}\}$ and
$\spn\{v_{42},v_{45}\}$.  The pair $(v_{12},v_{15})$ is an orthonormal basis
of the first subspace, while $(v_{13},v_{10})$ is an orthonormal basis of the
second.  Hence row~$1$ is an orthonormal basis.

For row~$2$, equation~\eqref{eq:completion} replaces the pair $(Q,S)$ in the
orthonormal basis $(e_0,P,R,Q,S)$ by the orthonormal pair $(v_{20},v_{21})$.
Similarly, in row~$3$ it replaces $(P,R)$ in the orthonormal basis
$(P,R,U,Q,H)$ by $(v_{30},v_{31})$.  Thus rows $2$ and $3$ are orthonormal.
For row~$4$, the pair $(Y,T)$ is perpendicular to
$(v_{42},v_{43},v_{45})$, and~\eqref{eq:v4041} replaces $(Y,T)$ by the
orthonormal pair $(v_{40},v_{41})$.  Finally,~\eqref{eq:v5051} replaces
$(v_{13},v_{43})$ by $(v_{50},v_{51})$.  It follows that
$(v_{50},v_{51},v_{53},Y,Z)$ is an orthonormal basis, which proves the claim
for row~$5$.

\emph{Columns $2,3,4,5$.}
Columns $2,3,4,5$ are respectively
\begin{align*}
&(e_1,v_{12},U,v_{42},Y),\\
&(e_2,v_{13},e_0,v_{43},v_{53}),\\
&(e_3,Y,P,Q,Z),\\
&(e_4,v_{15},R,H,v_{45}).
\end{align*}
Each is an orthonormal basis by the orthogonal decompositions used in the
construction.

\emph{Column $0$.}
This is the column for which the compatibility condition is needed.  All five
vectors in this column have unit norm.  The completion rule gives
$v_{10}\perp v_{20},v_{30}$, and
$\spn\{Q,S\}\perp\spn\{P,R\}$ gives $v_{20}\perp v_{30}$.  Moreover,
$v_{10}\perp v_{40},v_{50}$ because
$v_{10}\perp\spn\{Y,T,v_{13},v_{43}\}$, while
$v_{40}\perp v_{50}$ because
$\spn\{Y,T\}\perp\spn\{v_{13},v_{43}\}$.  It remains to verify the pairs
involving $v_{20}$ or $v_{30}$ with $v_{40}$ or $v_{50}$.

Put $q=\ip{v_{10}}{Q}$ and $p=\ip{v_{10}}{P}$.  Direct substitution, using
$N^2=1-\beta^2\varepsilon^2$, gives
\begin{align*}
\tau&=-\frac{\beta\varphi r_{43}}{N},&
\ip{v_{20}}{Y}&=-\frac{\beta\varphi q}{\rho_{20}},&
\ip{v_{20}}{T}&=-\frac{\alpha Nq}{r_{43}\rho_{20}},\\
\ip{v_{30}}{Y}&=-\frac{\beta\varepsilon p}{\rho_{30}},&
\ip{v_{30}}{T}&=-\frac{\kappa Np}{r_{43}\rho_{30}}.
\end{align*}
Consequently
\begin{align*}
\alpha\ip{v_{20}}{Y}+\tau\ip{v_{20}}{T}&=0,\\
\alpha\ip{v_{30}}{Y}+\tau\ip{v_{30}}{T}&=0,
\end{align*}
where the second equality uses $\varphi\kappa=\alpha\varepsilon$.  Since
$v_{40}=(\alpha Y+\tau T)/r_{40}$, this proves that both $v_{20}$ and
$v_{30}$ are perpendicular to $v_{40}$.  Finally, the definitions
in~\eqref{eq:aabb} and~\eqref{eq:v5051} give
\[
\ip{v_{20}}{v_{50}}=\frac{a_2a_1-a_1a_2}{r_{50}}=0,
\qquad
\ip{v_{30}}{v_{50}}=\frac{a_2b_1-a_1b_2}{r_{50}}=0,
\]
where the second equality uses~\eqref{eq:compatibility}.  Thus every pair of
distinct vectors in column~$0$ is orthogonal, so this column is an
orthonormal basis.

\emph{Column $1$.}
We have proved that every punctured row and each of the columns
$0,2,3,4,5$ is an orthonormal basis.  Lemma~\ref{lem:missing-column} now gives
the result for column~$1$.
\end{proof}

\end{document}
